% =============================================================================
% NOTA DO ARQUIVISTA — ARCHIVIST'S NOTE — 档案员按语
% Moldura diegética do Arquivo 1.260 (trilíngue)
% =============================================================================
\cleardoublepage
\thispagestyle{empty}
\chapter*{Nota do Arquivista --- Archivist's Note --- 档案员按语}
\addcontentsline{toc}{chapter}{Nota do Arquivista · Archivist's Note · 档案员按语}
\markboth{Nota do Arquivista · Archivist's Note · 档案员按语}{}

\begin{center}
\itshape\small
Este envelope contém a transcrição integral do diário do Paciente 1.260. A leitura do material original causou alterações psicossomáticas nos examinadores.\\[0.3em]
This envelope contains the complete transcription of the diary of Patient 1,260. Reading the original material caused psychosomatic alterations in the examiners.\\[0.3em]
此信封内装有患者1,260日记的完整誊本。阅读原始材料在检验人员身上引起了心身改变。\\[0.5em]
\upshape\scriptsize --- Anotação em caneta-tinteiro no envelope pericial nº 1.260/1979
\end{center}

\vspace{1.2em}

\begingroup

\noindent\textbf{Português.} O volume que você tem em mãos não foi escrito para ser publicado. É a reconstituição pericial do diário de \textbf{Joaquim Antônio Correia}, o Paciente 1.260 do Hospital Colônia de Barbacena. Sobrevivente da Seca de 1915 e dos campos de concentração de retirantes no Ceará, Joaquim foi internado no Colônia em 1917 e lá permaneceu por 62 anos, até sua morte em 1979. Ao longo desse cárcere, produziu 847 páginas de anotações fragmentárias. O caderno original foi encontrado amarrado ao seu peito por tiras de pano quando seu corpo foi exumado da vala comum --- 62 dias após o sepultamento. Segundo o relato pericial do coveiro, o papel emanava um odor adocicado de umidade e terra, e a tinta parecia não secar. Esta edição reúne 78 fragmentos organizados em cinco dimensões: as memórias da infância no sertão; o isolamento e as torturas no hospital psiquiátrico; a devassa dos laudos clínicos; a investigação forense conduzida pela Dra. Lúcia Mendes em 2026; e o registro da contaminação do próprio leitor. Não se trata de uma simples história de horror. A dor documentada nestas páginas não é ficção: ela atravessou décadas, consumiu médicos e peritos, e busca no ato da leitura um novo hospedeiro.

\vspace{0.8em}
\noindent{\itshape
Recomenda-se não ler este livro no escuro.\\[0.3em]
Recomenda-se não ler este livro em solidão absoluta.\\[0.3em]
Se o cheiro de terra molhada surgir durante a leitura, feche o volume imediatamente.}

\vspace{0.8em}
\noindent A ferida histórica sobre a qual esta narrativa se assenta é real. O Hospital Colônia de Barbacena (1903--1994) internou à revelia mais de 60.000 pessoas; 1.853 corpos foram exumados de suas valas comuns. Joaquim é um personagem, mas a dor que ele carrega --- fome, deslocamento, confinamento, tortura psiquiátrica --- pertenceu a milhares de brasileiros anônimos. O horror aqui evocado não deve ser consumido como entretenimento. A manipulação do leitor é uma escolha ética e formal --- não para enganar, mas para fazer o corpo sentir o que a mente sabe: que a vala não é uma metáfora, e que o sofrimento documentado existe para ser testemunhado, não estetizado.

\vspace{0.8em}
\noindent{\itshape Você será manipulado. Você será contaminado. Mas, ao final, saberá por quê.}

\vspace{1.2em}
\noindent\rule{\textwidth}{0.5pt}
\vspace{0.8em}

\noindent\textbf{English.} The volume you hold in your hands was not written to be published. It is the forensic reconstruction of the diary of \textbf{Joaquim Antônio Correia}, Patient 1,260 of the Barbacena Hospital-Colony. A survivor of the 1915 drought and of the concentration camps for refugees in Ceará, Joaquim was committed to the Colony in 1917 and remained there for 62 years, until his death in 1979. Throughout that imprisonment, he produced 847 pages of fragmentary notes. The original notebook was found tied to his chest with strips of cloth when his body was exhumed from the mass grave --- 62 days after burial. According to the gravedigger's forensic account, the paper gave off a sweet odor of dampness and earth, and the ink seemed never to dry. This edition brings together 78 fragments organized in five dimensions: the memories of childhood in the Sertão; the isolation and the tortures of the psychiatric hospital; the examination of the clinical reports; the forensic investigation conducted by Dr. Lúcia Mendes in 2026; and the record of the reader's own contamination. This is not a simple horror story. The pain documented in these pages is not fiction: it crossed decades, consumed doctors and examiners, and seeks, in the act of reading, a new host.

\vspace{0.8em}
\noindent{\itshape
It is recommended that you do not read this book in the dark.\\[0.3em]
It is recommended that you do not read this book in absolute solitude.\\[0.3em]
If the smell of wet earth arises during the reading, close the volume immediately.}

\vspace{0.8em}
\noindent The historical wound upon which this narrative rests is real. The Barbacena Hospital-Colony (1903--1994) interned more than 60,000 people without their consent; 1,853 bodies were exhumed from its mass graves. Joaquim is a character, but the pain he carries --- hunger, displacement, confinement, psychiatric torture --- belonged to thousands of anonymous Brazilians. The horror evoked here must not be consumed as entertainment. The manipulation of the reader is an ethical and formal choice --- not to deceive, but to make the body feel what the mind knows: that the trench is not a metaphor, and that the documented suffering exists to be witnessed, not aestheticized.

\vspace{0.8em}
\noindent{\itshape You will be manipulated. You will be contaminated. But in the end, you will know why.}

\vspace{1.2em}
\noindent\rule{\textwidth}{0.5pt}
\vspace{0.8em}

\noindent\textbf{中文.} 您手中这部书并非为出版而写。它是巴巴塞纳收容院患者1,260\textbf{若阿金·安东尼奥·科雷亚}日记的法证重建。1915年大旱与塞阿拉难民集中营的幸存者若阿金，于1917年被送入收容院，并在那里度过了62年，直到1979年去世。在那段监禁岁月中，他写下了847页零散笔记。遗体从乱葬坑中被掘出时——下葬62天后——原笔记本用布条绑在他的胸前。根据掘墓人的法证陈述，纸张散发出潮湿与泥土的甜味，墨迹似乎永远不会干。本版汇集了按五个维度组织的78个碎片：塞尔唐童年的记忆；精神病院的隔离与折磨；临床报告的检视；卢西亚·门德斯医生于2026年进行的法证调查；以及读者自身受污染的记录。这不是一个简单的恐怖故事。这些书页中记载的痛苦不是虚构：它跨越了数十年，消耗了医生和检验人员，并在阅读这一行为中寻找新的宿主。

\vspace{0.8em}
\noindent{\itshape
建议不要在黑暗中阅读本书。\\[0.3em]
建议不要完全独处时阅读本书。\\[0.3em]
若阅读中升起湿土的气味，请立即合上本书。}

\vspace{0.8em}
\noindent 这段叙事所依托的历史创伤是真实的。巴巴塞纳收容院（1903—1994）未经同意收容了六万多人；其乱葬坑中掘出1,853具遗体。若阿金是一个角色，但他所承载的痛苦——饥饿、流离、囚禁、精神科酷刑——属于成千上万无名的巴西人。此处唤起的恐怖不应作为娱乐消费。对读者的操纵是一种伦理与形式的选择——不是为了欺骗，而是为了让身体感受到心灵已知之事：乱葬坑不是隐喻，被记载的苦难存在是为了被见证，而不是被审美化。

\vspace{0.8em}
\noindent{\itshape 您将被操纵。您将被污染。但到头来，您会知道为什么。}

\vspace{1.4em}
\hfill{\scshape --- O Arquivista --- The Archivist --- 档案员}\hspace{1em}

\endgroup
\clearpage
